\begin{proof}
Since $C$ is finite, the cofinite-branch condition in the CLD topology is
automatic.  Hence \eqref{eq:p071-bivalent-unfolding} is a homeomorphism.
The image of a neighbourhood of $\eta_C$ contains an ordinary interval
about $0$, while neighbourhoods of the two full endpoints become
constant-collar neighbourhoods of $-\infty$ after using $q=\chi(y)$ or
$q=-\chi(y)$, respectively.

We first prove the easy inclusion.  Let $g$ be an exact two-star product
generator and unfold its second factor.  On the two half-spaces $q<0$ and
$q>0$, every first-factor branch function is smooth and jointly convex in
its branch coordinate and $q$.  Along $q=0$, the two restrictions have the
same value.  Their derivatives tangent to the first-factor branch agree
because they are derivatives of the common bivalent-node slice, and
\eqref{eq:p071-bivalent-second-balance} says that their one-sided normal
derivatives agree after the sign reversal in
\eqref{eq:p071-bivalent-unfolding}.  The unfolded branch function is
therefore continuously differentiable across $q=0$.  Its restriction to
every line segment is convex on the two pieces cut out by $q=0$ and has
matching one-sided derivatives at the cutting point.  It is consequently
jointly convex on the whole finite product chart.

At a node of $X_B$, the finite activity condition and
\eqref{eq:p071-bivalent-first-balance} give the star--Euclidean product
junction inequality by the derivative calculation in the proof of
\Cref{thm:p071-infinite-star-euclidean-collapse}.  This calculation remains
valid at $q=0$: the unfolded centre slice is continuously differentiable
there, and the branch functions are jointly convex across the same
hyperplane.  Exact terminal collars give the required traces at every full
endpoint.  Thus every generator belongs to $\mathcal J_{B,C}$.

The sheaf $\mathcal J_{B,C}$ is local by its definition.  It is closed
under pointwise decreasing sequential limits because, on each of its three
types of unfolded chart, it is one of the comparison sheaves proved
decreasing-sequentially closed in
\Cref{thm:p071-infinite-star-euclidean-collapse,%
thm:p071-infinite-star-infinite-collar-product-collapse}.  Minimality now
gives
\[
                         \mathcal P_{B,C}\subseteq\mathcal J_{B,C}.
\]

Conversely, let $u\in\mathcal J_{B,C}$ near $(x,y)$.  If
$q=\chi(y)$ is finite, choose an ordinary open interval $I\ni q$ in the
unfolded second factor.  This includes the case $q=0$.  The fixed-domain
assertion of
\Cref{thm:p071-infinite-star-euclidean-collapse}, with $d=1$, gives a
pointwise decreasing sequence of exact star--Euclidean generators on one
fixed smaller product neighbourhood in $X_B\times I$.  Pulling them back
by $\chi$ gives members of $\mathcal G_{B,C,0}$.  At $q=0$ they are smooth
across the unfolded coordinate, so
\eqref{eq:p071-bivalent-second-balance} holds; this is stronger than the
piecewise smoothness required of a two-star generator.

If $q=-\infty$, apply instead the fixed-domain assertion of
\Cref{thm:p071-infinite-star-infinite-collar-product-collapse} and pull the
approximants back by $\chi$.  If $q=+\infty$, first reflect the unfolded
coordinate and use the same theorem.  In both cases the approximants have
the exact terminal collar in the second factor, in addition to the exact
first-factor collars and balance law.

We have proved $\mathcal J_{B,C}\subseteq\mathcal P_{B,C,1}$ at every
point.  Together with the tautological inclusion
$\mathcal P_{B,C,1}\subseteq\mathcal P_{B,C}$ and the easy inclusion above,
this proves \eqref{eq:p071-bivalent-product-collapse} and the
fixed-neighbourhood assertion.  Interchanging the two factors proves the
final symmetry statement.
\end{proof}
